\chapter{Implementation and Testing}
\label{sec:implementation}

This chapter documents the technical implementation details, algorithms, external dependencies, and testing strategies for the personalized-ai-storybook FYP-2 system. The implementation demonstrates practical application of AI technologies across six new modules (TTS, STT, styled PDF, evaluation, personalized story, and UI enhancements).

\section{Algorithm Design}

\subsection{Multi-Agent Story Generation Pipeline (Enhanced)}

The core story generation algorithm remains a sequential multi-agent pipeline, extended with background evaluation, TTS triggering, and STT preprocessing.

\begin{verbatim}
ALGORITHM: Enhanced Multi-Agent Story Generation
INPUT:  user_prompt (string | audio), generate_images (bool),
        use_personalized_images (bool), user_photo (image | None),
        voice (string), genre, style, length
OUTPUT: complete_story (JSON), pdf_path (string), status

BEGIN
  1. IF audio_input THEN
       audio = STTManager.transcribe_audio(audio_file)
       user_prompt = audio.text
  
  2. PromptAgent.process_prompt(user_prompt)
     - Validate length [5-500 chars]
     - Classify: normal | detailed | invalid | nonsense
     - Apply ContentSafety filter (blocked keywords)
     - Enhance simple prompts with genre/style context
     - RETURN enhanced_prompt, prompt_type

  3. IF prompt_type IN {invalid, nonsense}
       RETURN error_response("Invalid prompt")

  4. DirectorAgent.create_story(enhanced_prompt)
     4.1 WriterAgent.generate_story(enhanced_prompt)
         - Format LLM prompt with genre/style/length metadata
         - Call Ollama (llama3.1:8b-instruct-q4_K_M)
         - Parse JSON response; retry up to 2 times
         - RETURN story_draft
     4.2 ReviewerAgent.review_story(story_draft)
         - Validate coherence, age-appropriateness, scene count
         - RETURN reviewed_story
     4.3 EditorAgent.edit_story(reviewed_story)
         - Finalize metadata and formatting
         - export_pdf(story_data) -> styled PDF
         - RETURN final_story

  5. IF use_personalized_images AND user_photo THEN
       OllamaManager.pause_ollama()  // Free GPU VRAM
       FOR each scene IN final_story.scenes DO
         PersonalizedImageAgent.generate_image(
           scene.description, user_photo, ip_adapter_weights)
       OllamaManager.resume_ollama()
     ELSE IF generate_images THEN
       OllamaManager.pause_ollama()
       FOR each scene IN final_story.scenes DO
         ImageAgent.generate_image(scene.description)
       OllamaManager.resume_ollama()

  6. Save final_story to generated/stories/
  7. story_json["mode"] = ("personalized" | "simple")
  8. Trigger EvaluationAgent (background subprocess)
  9. RETURN final_story, pdf_path, status
END
\end{verbatim}

\subsection{TTS Audio Generation with Caching}

The TTS algorithm ensures efficient narration through a cache-first approach.

\begin{verbatim}
ALGORITHM: TTS Audio Generation (Piper, Offline)
INPUT:  text (string), story_id (int), scene_number (int),
        voice ("male" | "female")
OUTPUT: relative_audio_path (string)

BEGIN
  1. voice_model = VOICES[voice]
     // female: en_US-lessac-medium
     // male:   en_US-ryan-medium
  
  2. filename = f"story_{id}_scene_{n}_{voice}.wav"
     output_path = join(AUDIO_DIR, filename)
  
  3. IF exists(output_path) AND size > 0 THEN
       RETURN "/generated/audio/" + filename  // Cache hit

  4. Write text to temporary .txt file
  5. Build PowerShell command:
       "Get-Content temp.txt | python -m piper
        --model model.onnx
        --config model.onnx.json
        --output_file output.wav"
  
  6. subprocess.run(cmd, shell=True)
  7. Delete temp .txt file
  8. IF returncode != 0 OR file empty THEN
       RETURN ""  // Graceful failure
  
  9. RETURN "/generated/audio/" + filename
END
\end{verbatim}

\subsection{STT Speech Transcription (ffmpeg-Free)}

The STT algorithm uses only Python built-ins for WAV processing, eliminating the ffmpeg system dependency.

\begin{verbatim}
ALGORITHM: Speech-to-Text Transcription (Whisper, Offline)
INPUT:  audio_file_path (string), language (string)
OUTPUT: {text: string, language: string}

BEGIN
  1. model = WhisperSTTManager.load_model()
     // Lazy loads "base" Whisper model on first call

  2. raw_bytes = read_binary(audio_file_path)
  
  3. // WAV decoding (no ffmpeg)
     wave.open(BytesIO(raw_bytes)) ->
       n_channels, sample_width, sample_rate, n_frames, raw
     
     IF sample_width == 2 THEN dtype = int16, max = 32768
     IF sample_width == 4 THEN dtype = int32, max = 2147483648
     samples = frombuffer(raw, dtype).astype(float32) / max

  4. IF n_channels > 1 THEN
       samples = samples.reshape(-1, n_channels).mean(axis=1)
       // Stereo -> Mono

  5. // Resample to 16kHz (Whisper requirement)
     target_len = int(len(samples) / orig_sr * 16000)
     audio_16k = interp(
       linspace(0, len(samples), target_len),
       arange(len(samples)), samples).astype(float32)

  6. result = model.transcribe(audio_16k, language, fp16=False)
  7. RETURN {text: result["text"].strip(),
             language: result["language"]}
END
\end{verbatim}

\subsection{Evaluation Metrics Algorithm}

The evaluation algorithm computes six quality metrics across multiple ML models.

\begin{verbatim}
ALGORITHM: Comprehensive Story Evaluation
INPUT:  story_id, story_dict, image_paths, mode
OUTPUT: evaluation_results (JSON with 0-100 scores)

BEGIN
  1. texts = [scene.text for scene in story.scenes]
  2. clip_texts = [scene.image_description for scene in scenes]
  3. full_text = join(texts)

  // Metric 1: Image-Text Alignment (CLIP)
  4. FOR each (img, text) IN zip(image_paths, clip_texts):
       inputs = CLIPProcessor(text, image)
       outputs = CLIPModel(inputs)
       cos_sim = cosine_similarity(image_embeds, text_embeds)
     image_text_score = scale([0.15, 0.40] -> [60, 100])

  // Metric 2: Visual Consistency (CLIP)
  5. embeddings = [CLIPModel.image_features(img) for img]
     pairwise_sims = cosine_similarity(embeddings_matrix)
     visual_score = scale([0.50, 1.0] -> [60, 100])

  // Metric 3: Character Consistency (InsightFace, if personalized)
  6. IF mode == "personalized":
       face_embeddings = [InsightFace.detect(img)[0].embedding]
       face_similarities = [dot(e1, e2) for consecutive pairs]
       character_score = scale([0.20, 0.80] -> [60, 100])

  // Metric 4: Text Coherence (Sentence Transformers)
  7. embeddings = SentenceTransformer.encode(texts)
     consecutive_sims = [cosine_sim(e[i], e[i+1])]
     coherence_score = scale([0.25, 0.85] -> [65, 100])

  // Metric 5: Readability (Flesch-Kincaid)
  8. fk_grade = textstat.flesch_kincaid_grade(full_text)
     // Grade 3-4: 95, Grade 5-6: 90, Grade 7-8: 78, etc.

  // Metric 6: Story Structure
  9. IF first_scene.len > 50: total += 34
     IF num_scenes >= 3: total += 33
     IF last_scene.len > 50: total += 33

  10. weights = {image_text: 0.20, visual: 0.15,
                 coherence: 0.25, readability: 0.20,
                 structure: 0.15}
  11. overall = weighted_average(scores, weights)

  12. results = {story_id, mode, overall_score,
                 metrics: {each_metric: {score, label}},
                 evaluation_time_seconds}
  13. save_json(results, "generated/evaluations/")
  14. RETURN results
END
\end{verbatim}

\subsection{Styled Open-Book PDF Generation}

The PDF generation algorithm recreates the UI open-book aesthetic in a downloadable document.

\begin{verbatim}
ALGORITHM: Styled Open-Book PDF Export (ReportLab)
INPUT:  story_data (JSON), output_filename (string)
OUTPUT: pdf_path (string)

BEGIN
  1. Register Fonts:
     Try load Georgia (.ttf) from C:\Windows\Fonts\
     Fallback: Times-Roman (built-in ReportLab)
  
  2. Compute geometry:
     PAGE_W, PAGE_H = A4 landscape (841.89 x 595.28 pts)
     LEATHER_PAD = 12 pts, SPINE_WIDTH = 8 pts
     LEFT_PAGE = [12, 12, half-width-8, height-24]
     RIGHT_PAGE = [half-width+8, 12, half-width-8, height-24]

  3. Cover Page:
     - Fill background with C_LEATHER (#4A352C)
     - Draw cream left page with last scene image (white mat border)
     - Draw deep purple right page (C_COVER_PLM #23102E)
     - Render story title in gold (C_GOLD_LGT) font,
       dynamically wrapped; decorative gold lines above/below
     - Draw spine crease (gradient shadows + spine tube)
     - Add "Cover Page" pill indicator

  4. FOR each scene IN scenes:
     - Draw leather background + two cream pages + spine
     LEFT PAGE (illustration):
       - Scale image to fit within padded area
       - Draw drop shadow offset
       - Draw white photo mat (12pt padding)
       - Render image centered
     RIGHT PAGE (text):
       - Draw red margin line (C_RED_MARG)
       - AUTO-SCALE font [14pt, 13, 12, 11, 10pt]:
           wrap text to content_width
           test if total_height <= content_height
           break when fits or at 10pt
       - Draw notebook rulings (C_NOTE_LINE) aligned with text
       - Render wrapped text lines on rulings (first line indented)
       - Draw "Page X of N" pill at bottom center
     c.showPage()

  5. canvas.save()
  6. RETURN pdf_path
END
\end{verbatim}


\section{External APIs/SDKs}

\begin{table}[H]
\centering
\caption{External APIs/SDKs Used in FYP-2 Implementation}
\resizebox{\textwidth}{!}{%
\begin{tabular}{|p{0.22\textwidth}|p{0.3\textwidth}|p{0.28\textwidth}|p{0.20\textwidth}|}
\hline
\textbf{API/SDK} & \textbf{Description} & \textbf{Purpose of Usage} & \textbf{Key Classes/Endpoints} \\ \hline
Ollama (Latest) & Local LLM server for running quantized models & Story text generation, workshop chat & \texttt{ollama run llama3.1:8b} \\ \hline
stable-diffusion-webui & AUTOMATIC1111 WebUI with local REST API & IP-Adapter integration, personalized/simple image generation & \texttt{/sdapi/v1/txt2img}, \texttt{/sdapi/v1/img2img} \\ \hline
Piper TTS & Offline ONNX TTS engine, run via Python \texttt{-m piper} & Scene-by-scene voice narration, offline WAV generation & CLI subprocess via PowerShell \\ \hline
OpenAI Whisper & Local STT model for offline transcription & Voice story prompt submission & \texttt{whisper.load\_model()}, \texttt{model.transcribe()} \\ \hline
CLIP (ViT-B/32) & Contrastive Language-Image Pretraining & Image-text alignment \& visual consistency evaluation & \texttt{CLIPModel}, \texttt{CLIPProcessor} \\ \hline
Sentence Transformers & Dense sentence embedding model & Text coherence evaluation between scenes & \texttt{SentenceTransformer}, \texttt{encode()} \\ \hline
InsightFace & Face recognition and embedding library & Character face consistency evaluation (personalized mode) & \texttt{FaceAnalysis}, \texttt{.get(image)} \\ \hline
FastAPI 0.100+ & Async Python web framework & Backend REST API server & \texttt{@app.post()}, \texttt{UploadFile} \\ \hline
Diffusers 0.21+ & Hugging Face diffusion model library & Fallback image generation pipeline & \texttt{StableDiffusionPipeline} \\ \hline
ReportLab 4.0+ & PDF generation library & Styled open-book PDF export & \texttt{canvas.Canvas}, \texttt{roundRect}, \texttt{drawImage} \\ \hline
PyTorch 2.0+ & Deep learning framework & Model inference, tensor operations, CUDA/CPU device management & \texttt{torch.cuda}, \texttt{torch.nn.functional.cosine\_similarity} \\ \hline
SQLAlchemy 2.0+ & Python ORM for PostgreSQL & Database CRUD operations & \texttt{Session}, \texttt{Base.metadata.create\_all()} \\ \hline
Next.js 14+ & React framework with SSR & Frontend application, open-book UI & \texttt{pages/}, \texttt{app/} routing \\ \hline
LangChain & AI workflow orchestration & Idea Workshop Agent multi-turn conversation & \texttt{LLMChain}, memory buffers \\ \hline
textstat & Python readability metrics & Flesch-Kincaid grade evaluation & \texttt{textstat.flesch\_kincaid\_grade()} \\ \hline
PIL (Pillow) & Python Imaging Library & Image preprocessing for PDF, IP-Adapter & \texttt{Image.open()}, \texttt{.convert("RGB")} \\ \hline
\end{tabular}%
}
\label{tab:external_apis}
\end{table}


\section{Testing Details}

The system employs a comprehensive testing strategy combining black-box and white-box approaches across 103 total test cases.

\subsection{White-Box (Unit) Testing}

White-box tests were implemented using \textbf{PyTest} with \textbf{Coverage.py}, targeting critical business logic modules.

\textbf{Testing Summary}: 35 unit tests, 35 passed, 0 failed (100\% pass rate).

\begin{table}[H]
\centering
\caption{Unit Test Coverage by Module}
\begin{tabular}{|l|c|c|c|c|}
\hline
\textbf{Module} & \textbf{Statements} & \textbf{Covered} & \textbf{Coverage} & \textbf{Status} \\ \hline
auth/schemas.py & 37 & 37 & 100.00\% & Excellent \\ \hline
utils/content\_safety.py & 47 & 46 & 97.87\% & Excellent \\ \hline
agents/prompt\_agent.py & 30 & 26 & 86.67\% & Very Good \\ \hline
agents/story\_agent.py & 13 & 12 & 92.31\% & Excellent \\ \hline
agents/chatbot\_agent.py & 53 & 39 & 73.58\% & Good \\ \hline
\multicolumn{3}{|l|}{\textbf{Tested Modules Average}} & \textbf{90.09\%} & \\ \hline
\end{tabular}
\label{tab:coverage}
\end{table}

\textbf{Test Case Samples}:

\textit{Prompt Agent Tests (FR1.1 -- FR1.3)}:
\begin{itemize}
    \item \texttt{test\_empty\_input\_handling()}: Input: empty string. Expected: classified as invalid. \textbf{Result: PASS}.
    \item \texttt{test\_nonsense\_rejection()}: Input: ``asdfghjkl''. Expected: HTTP 400 error. \textbf{Result: PASS}.
    \item \texttt{test\_short\_prompt\_classification()}: Input: ``cat''. Expected: classified as ``short''. \textbf{Result: PASS}.
    \item \texttt{test\_url\_rejection()}: Input: ``https://example.com''. Expected: rejected. \textbf{Result: PASS}.
\end{itemize}

\textit{Content Safety Tests (FR1.11)}:
\begin{itemize}
    \item \texttt{test\_block\_explicit\_content()}: Input: prompt with explicit keyword. Expected: blocked. \textbf{Result: PASS}.
    \item \texttt{test\_allow\_safe\_content()}: Input: ``A brave fox on an adventure''. Expected: passed. \textbf{Result: PASS}.
    \item \texttt{test\_case\_insensitive\_blocking()}: Input: ``VIOLENCE''. Expected: blocked regardless of case. \textbf{Result: PASS}.
\end{itemize}

\textit{Story Agent Tests (FR1.4)}:
\begin{itemize}
    \item \texttt{test\_agent\_initialization()}: Verifies DirectorAgent composed correctly. \textbf{Result: PASS}.
    \item \texttt{test\_story\_generation\_via\_director()}: Verifies pipeline execution and JSON output. \textbf{Result: PASS}.
\end{itemize}

\textit{Authentication Schema Tests}:
\begin{itemize}
    \item \texttt{test\_valid\_user\_creation()}: Valid username, password, email. \textbf{Result: PASS}.
    \item \texttt{test\_password\_too\_short()}: Password = ``ab''. Expected: ValidationError. \textbf{Result: PASS}.
    \item \texttt{test\_email\_invalid\_format()}: Email = ``notanemail''. Expected: ValidationError. \textbf{Result: PASS}.
\end{itemize}

\subsection{Black-Box Testing}

Black-box testing was conducted through documented test cases without knowledge of internal implementation, testing system behavior from the user perspective. Total: \textbf{68 test cases}.

\begin{table}[H]
\centering
\caption{Black-Box Testing Summary}
\begin{tabular}{|l|c|p{7cm}|}
\hline
\textbf{Technique} & \textbf{Test Cases} & \textbf{Coverage Areas} \\ \hline
Equivalence Class Partitioning & 30 & Prompt length classes, valid/invalid modes, voice types, image formats, genre/style values \\ \hline
Boundary Value Analysis & 26 & Prompt at exactly 5 chars (min), 500 chars (max), 1 scene vs 3 scenes, empty audio file, minimum image dimensions \\ \hline
Functional Test Cases & 12 & End-to-end story generation, TTS playback, STT transcription, PDF export, character chatbot, idea workshop, evaluation display \\ \hline
\textbf{Total} & \textbf{68} & \\ \hline
\end{tabular}
\label{tab:blackbox}
\end{table}

\textbf{Selected Functional Test Cases}:

\textit{TC-F01 -- End-to-End Story Generation}:
\begin{itemize}
    \item Input: Prompt ``A little dragon learns to fly'', genre=Fantasy, style=Cartoon, images=enabled.
    \item Steps: POST /api/generate → verify response JSON → verify images in generated/images/ → verify PDF in generated/exports/.
    \item Expected: Story with 3 scenes, 3 images, PDF exported. \textbf{Result: PASS}. Average time: 2.5 minutes (GPU).
\end{itemize}

\textit{TC-F02 -- TTS Audio Generation and Caching}:
\begin{itemize}
    \item Input: POST /api/tts \{story\_id: 1, scene\_number: 1, voice: "female"\}.
    \item First call: WAV generated; Second call: cached file returned immediately.
    \item Expected: Audio file at /generated/audio/story\_1\_scene\_1\_female.wav. \textbf{Result: PASS}.
\end{itemize}

\textit{TC-F03 -- STT Voice Transcription}:
\begin{itemize}
    \item Input: WAV file with spoken sentence ``A brave fox on a magical adventure''.
    \item Expected: Transcribed text closely matches input. \textbf{Result: PASS}. Transcription accuracy: $\sim$95\%.
\end{itemize}

\textit{TC-F04 -- Styled PDF Export}:
\begin{itemize}
    \item Input: Generated story with 3 scenes and images.
    \item Expected: PDF with leather cover, 3 scene spreads, Georgia font. \textbf{Result: PASS}. Export time: 3.2s.
\end{itemize}

\textit{TC-F05 -- Personalized Story Generation}:
\begin{itemize}
    \item Input: Prompt + reference photo (JPEG, 512$\times$512).
    \item Expected: 3 illustrations with consistent facial likeness across all scenes. \textbf{Result: PASS}.
\end{itemize}

\textit{TC-F06 -- Evaluation Metrics}:
\begin{itemize}
    \item Input: Completed simple story with 3 images.
    \item Expected: JSON report with image\_text\_score, visual\_consistency, text\_coherence, readability, story\_structure, overall\_score. \textbf{Result: PASS}. Average overall score: $\sim$82/100.
\end{itemize}


\subsection{Performance Testing}

\begin{table}[H]
\centering
\caption{System Performance Benchmarks (NVIDIA GTX 1060, Intel i5, 16GB RAM)}
\begin{tabular}{|l|c|c|}
\hline
\textbf{Operation} & \textbf{GPU System} & \textbf{CPU-Only} \\ \hline
Story text generation (3 scenes) & 12--15 seconds & 18--22 seconds \\ \hline
Single image generation & 35--45 seconds & 2--3 minutes \\ \hline
Complete 3-scene story (with images) & 2--2.5 minutes & 8--10 minutes \\ \hline
Styled PDF export & 3--5 seconds & 3--5 seconds \\ \hline
TTS per scene (Piper, first gen) & 2--4 seconds & 2--4 seconds \\ \hline
TTS per scene (cached) & $<$100ms & $<$100ms \\ \hline
STT transcription (15s audio) & 3--6 seconds & 5--8 seconds \\ \hline
Evaluation (3 scenes, CLIP + all) & 15--30 seconds & 30--60 seconds \\ \hline
Database query latency & $<$100ms & $<$100ms \\ \hline
\end{tabular}
\label{tab:performance}
\end{table}
